\section{Security}

\subsection{Who is calling? Authentication}

\begin{description}
\item[Students and admins] sign in through the university's identity provider
  using \textbf{OpenID Connect} (OIDC): the browser is sent to the campus login
  page, the user signs in there, and the identity provider sends back a signed
  token saying who they are. The platform never sees passwords. For
  development a mock identity provider (\texttt{oidc\_dev.py}) accepts any
  e-mail, and an \texttt{X-Dev-User} header skips login entirely.
\item[Lab PCs] receive a random 32-byte token when they first register. The
  server stores only a \emph{hash} of it:
\end{description}

\begin{lstlisting}[language=Python]
# control/app/crypto.py
def hash_token(token: str) -> str:
    return hashlib.sha256(f"{settings.token_pepper}:{token}".encode()).hexdigest()
\end{lstlisting}

If the database leaks, the attacker gets hashes, not usable tokens. The
\emph{pepper} is a server-side secret mixed into every hash, so the hashes are
useless without it too. Revoking a stolen PC's token
(\texttt{POST /nodes/<id>/revoke}) sets \texttt{revoked\_at}, and the
authentication dependency rejects it from then on.

\subsection{Refusing to run unsafely}

Development defaults are convenient and dangerous: the dev header lets anyone
be anyone. Rather than trusting a checklist, the API checks itself at startup
when \texttt{ENVIRONMENT=production}:

\begin{lstlisting}[language=Python]
def production_problems(self) -> list[str]:
    problems = []
    if self.dev_auth:
        problems.append("DEV_AUTH must be 0 (the X-Dev-User header ...)")
    if len(self.token_pepper) < 24:
        problems.append("TOKEN_PEPPER must be a random secret of 24+ chars")
    if self.minio_secret_key in {"labminio", ""}:
        problems.append("MINIO_SECRET_KEY is the dev default")
    ...
    return problems
\end{lstlisting}

With any problem present, startup fails with a list of exactly what to fix. A
misconfigured server never comes up half-safe. This was tested by starting the
production container with \texttt{DEV\_AUTH=1} and a short pepper; it refused
and named both.

\subsection{HTTPS with Caddy}

In production, Caddy sits in front of everything as a reverse proxy and
terminates TLS for three hostnames (portal, API, registry). Two certificate
options:

\begin{itemize}
\item \textbf{ACME} (Let's Encrypt): real, publicly trusted certificates,
  when the hostnames are reachable from the internet.
\item \textbf{Internal CA} (\texttt{tls internal}): Caddy runs its own
  certificate authority. Machines must be told to trust its root certificate
  once; the lab PC installer's \verb|--ca-cert| flag does that for both the
  agent and Docker.
\end{itemize}

\subsection{Read for everyone, write for one}

Every PC must pull images; only the builder should push them. Caddy enforces
that with HTTP methods alone:

\begin{verbatim}
{$REGISTRY_HOST} {
    @write not method GET HEAD
    basic_auth @write {
        {$REGISTRY_PUSH_USER} {$REGISTRY_PUSH_HASH}
    }
    reverse_proxy registry:5000
}
\end{verbatim}

Tested end to end: an anonymous pull returned 200, an anonymous push 401, and
the builder's push 202.

\subsection{Least privilege on the PC}

The agent's Linux user may start the Docker daemon at boot, and nothing else as
root. The installer adds exactly one sudo rule and verifies it with
\texttt{visudo}:

\begin{verbatim}
lab ALL=(root) NOPASSWD: /usr/sbin/service docker start
\end{verbatim}

A test in a fresh Ubuntu container confirmed \texttt{sudo service docker start}
works for that user and \texttt{sudo id} is refused.

\subsection{Thinking like an attacker}

A short threat model for this system:

\begin{center}
\begin{tabular}{ll}
\textbf{Threat} & \textbf{Defence} \\
Student script reads others' data & container, read-only input, no network \\
Script attacks the PC & non-root user, no docker.sock, cgroup limits \\
Fake agent update over Wi-Fi & RSA signature, pinned public key \\
Stolen PC replays its token & token revocation, hashed tokens \\
Student impersonates another & OIDC login, dev header off in production \\
Anyone pushes a malicious image & registry push requires credentials \\
Database leak & only peppered token hashes stored \\
\end{tabular}
\end{center}

Writing a table like this, even a rough one, is the first step of any security
review (the formal versions are called STRIDE or attack trees).
